\subsection{因果狄拉克场 }
一般形式的场是：
\begin{equation}
    \Psi_{ab}(x) =(\frac{1}{2\pi})^\frac{3}{2}\sum_{\sigma}\int d^3p \, \left\{(-)^{2B} u_{ab}(\boldsymbol{q},\sigma)b^\dagger(\boldsymbol{q},\sigma) e^{-ipx} + v_{ab}(\boldsymbol{q},\sigma) a(\boldsymbol{q},\sigma)e^{ipx} \right\}
\end{equation}
上述我们知道静止系系数 $u(0)$ 和 $v(0)$ 是正比于Clebsch-Gordan  系数 
\begin{equation}
    u_{ab}(0, \sigma) = \frac{1}{\sqrt{2m}} \langle A, a; B, b | j, \sigma \rangle
    \label{eq:u0_general_1/2}
\end{equation}
通过 Boost 获得有限动量系数 $u(\boldsymbol{q})$ 和 $v(\boldsymbol{q})$的一般公式是：
\begin{equation}
    u_{ab}(\boldsymbol{q}, \sigma) = \frac{1}{\sqrt{2q^0}} \sum_{\bar{a},\bar{b}} \left[e^{-\boldsymbol{\phi}\cdot\mathbf{J}^{(A)}}\right]_{a\bar{a}} \left[e^{+\boldsymbol{\phi}\cdot\mathbf{J}^{(B)}}\right]_{b\bar{b}} \langle A,\bar{a}; B,\bar{b} | j, \sigma \rangle
    \label{eq:u_finite_general_1/2}
\end{equation}
其中$\boldsymbol{\phi} = \hat{\boldsymbol{q}} \text{arctanh}(|\boldsymbol{q}|/q^0)$
\vspace{0.5cm}
\begin{center}
    \large\text{$(A, B) = (1/2, 0)\oplus(1/2, 0)$}
\end{center}
\vspace{0.5cm}
对于自旋 $j=1/2$ 的有质量粒子，在静止参考系（$\boldsymbol{p}=\mathbf{0}$，$p^0=m$）中，狄拉克方程 $(\gamma^\mu p_\mu - m)u=0$ 退化为：
\begin{equation}
    (m\gamma^0 - m)u(0,\sigma) = 0 \quad\Longrightarrow\quad \gamma^0 u(0,\sigma) = u(0,\sigma)
\end{equation}
这是 $\gamma^0$ 矩阵的本征值为 $+1$ 的本征方程。在手征表示下 $\gamma^0 = \begin{pmatrix} 0&I\\I&0 \end{pmatrix}$，写出分量形式：
\begin{equation}
    \begin{pmatrix} 0&I\\I&0 \end{pmatrix}\begin{pmatrix} u_R \\ u_L \end{pmatrix} = \begin{pmatrix} u_R \\ u_L \end{pmatrix} \quad\Longrightarrow\quad u_L = u_R
\end{equation}
因此本征值 $+1$ 的本征空间是二维的（$u_R = u_L$ 可取两个线性独立基矢），对应两个自旋投影 $\sigma = \pm 1/2$。

选择洛伦兹标量归一化 $\bar{u}u = u^\dagger\gamma^0 u = 2m$，逐一写出具体的四分量列向量：

\textbf{$\sigma = +1/2$：}
\begin{equation}
    u(0, +\tfrac{1}{2}) = \sqrt{m}\begin{pmatrix} 1\\0\\1\\0 \end{pmatrix} = \sqrt{m}\begin{pmatrix} \xi_{+1/2} \\ \xi_{+1/2} \end{pmatrix}
\end{equation}

\textbf{$\sigma = -1/2$：}
\begin{equation}
    u(0, -\tfrac{1}{2}) = \sqrt{m}\begin{pmatrix} 0\\1\\0\\1 \end{pmatrix} = \sqrt{m}\begin{pmatrix} \xi_{-1/2} \\ \xi_{-1/2} \end{pmatrix}
\end{equation}
其中 $\xi_{+1/2} = \binom{1}{0}$，$\xi_{-1/2} = \binom{0}{1}$。简记为 $u(0,\sigma) = \sqrt{m}\binom{\xi_\sigma}{\xi_\sigma}$。

\textbf{$\sigma = +1/2$：} 
\begin{equation}
    (-)^{1}(-)^{1+1/2}u(0, +\tfrac{1}{2}) = \sqrt{m}\begin{pmatrix} 0\\1\\0\\-1 \end{pmatrix} = \sqrt{m}\begin{pmatrix} \chi_{+1/2} \\ -\chi_{+1/2} \end{pmatrix}
\end{equation}

\textbf{$\sigma = -1/2$：} 
\begin{equation}
        (-)^{1}(-)^{1+1/2}u(0, -\tfrac{1}{2})= \sqrt{m}\begin{pmatrix} -1\\0\\1\\0 \end{pmatrix} = \sqrt{m}\begin{pmatrix} \chi_{-1/2} \\ -\chi_{-1/2} \end{pmatrix}
\end{equation}
简记为 $v(0,\sigma) = \sqrt{m}\binom{\chi_\sigma}{-\chi_\sigma}$。
\vspace{2ex}\par
在直和表示 $(1/2,0)\oplus(0,1/2)$ 中，Boost 矩阵是对角分块的（上块按右手加速，下块按左手加速）：
\begin{equation}
    S(L(p)) = \begin{pmatrix} e^{+\frac{1}{2}\boldsymbol{\phi}\cdot\boldsymbol{\sigma}} & 0 \\ 0 & e^{-\frac{1}{2}\boldsymbol{\phi}\cdot\boldsymbol{\sigma}} \end{pmatrix}
\end{equation}
利用泡利矩阵的性质 $(\hat{\boldsymbol{p}}\cdot\boldsymbol{\sigma})^2 = I$ 以及双曲半角公式（$\cosh\phi = p^0/m$，$\sinh\phi = |\boldsymbol{p}|/m$）：
\begin{align}
    e^{+\frac{1}{2}\boldsymbol{\phi}\cdot\boldsymbol{\sigma}} &= \cosh\frac{\phi}{2}\cdot I + \sinh\frac{\phi}{2}\cdot\hat{\boldsymbol{p}}\cdot\boldsymbol{\sigma} = \frac{(p^0+m)I + \boldsymbol{p}\cdot\boldsymbol{\sigma}}{\sqrt{2m(p^0+m)}} \\
    e^{-\frac{1}{2}\boldsymbol{\phi}\cdot\boldsymbol{\sigma}} &= \cosh\frac{\phi}{2}\cdot I - \sinh\frac{\phi}{2}\cdot\hat{\boldsymbol{p}}\cdot\boldsymbol{\sigma} = \frac{(p^0+m)I - \boldsymbol{p}\cdot\boldsymbol{\sigma}}{\sqrt{2m(p^0+m)}}
\end{align}
将 Boost 矩阵作用于静止系旋量 $u(0,\sigma) = \sqrt{m}\binom{\xi_\sigma}{\xi_\sigma}$：
\begin{equation}
    u(\boldsymbol{p},\sigma) = S(L(p))\, u(0,\sigma) = \sqrt{m}\begin{pmatrix} e^{+\frac{1}{2}\boldsymbol{\phi}\cdot\boldsymbol{\sigma}}\xi_\sigma \\ e^{-\frac{1}{2}\boldsymbol{\phi}\cdot\boldsymbol{\sigma}}\xi_\sigma \end{pmatrix}
\end{equation}
四维记号下： $\sqrt{p\cdot\bar{\sigma}} \equiv e^{+\frac{1}{2}\boldsymbol{\phi}\cdot\boldsymbol{\sigma}}\sqrt{m}$ 和 $\sqrt{p\cdot\sigma} \equiv e^{-\frac{1}{2}\boldsymbol{\phi}\cdot\boldsymbol{\sigma}}\sqrt{m}$（其中 $p\cdot\sigma = p^0 I - \boldsymbol{p}\cdot\boldsymbol{\sigma}$，$p\cdot\bar{\sigma} = p^0 I + \boldsymbol{p}\cdot\boldsymbol{\sigma}$），则：
\begin{equation}
    \boxed{u(\boldsymbol{p},\sigma) = \begin{pmatrix} \sqrt{p\cdot\bar{\sigma}}\,\xi_\sigma \\ \sqrt{p\cdot\sigma}\,\xi_\sigma \end{pmatrix}}
\end{equation}
将 Boost 矩阵作用于 $v(0,\sigma) = \sqrt{m}\binom{\chi_\sigma}{-\chi_\sigma}$（注意下半部分的负号在 Boost 后被保留）：
\begin{equation}
    \boxed{v(\boldsymbol{p},\sigma) = \begin{pmatrix} \sqrt{p\cdot\bar{\sigma}}\,\chi_\sigma \\ -\sqrt{p\cdot\sigma}\,\chi_\sigma \end{pmatrix}}
\end{equation}
这些波函数在任意惯性系下自动满足 $(\slashed{p}-m)u=0$ 和 $(\slashed{p}+m)v=0$。
\vspace{2ex}\par
为了得到洛伦兹协变的结果，我们计算 $u\bar{u}$（其中 $\bar{u}=u^\dagger\gamma^0$）而非 $uu^\dagger$。

对 $\sigma = \pm 1/2$ 求和，利用 $\sum_\sigma\xi_\sigma\xi_\sigma^\dagger = I$，逐块计算 $4\times 4$ 矩阵的四个 $2\times 2$ 子块：
\begin{itemize}
    \item 左上块：$\sqrt{p\cdot\bar{\sigma}}\cdot I\cdot\sqrt{p\cdot\sigma} = \sqrt{(p\cdot\bar{\sigma})(p\cdot\sigma)}$。利用 $(p\cdot\bar{\sigma})(p\cdot\sigma) = (p^0)^2 - |\boldsymbol{p}|^2 = m^2$，得 $= mI$。
    \item 右上块：$\sqrt{p\cdot\bar{\sigma}}\cdot I\cdot\sqrt{p\cdot\bar{\sigma}} = p\cdot\bar{\sigma} = p^0 I + \boldsymbol{p}\cdot\boldsymbol{\sigma}$。
    \item 左下块：$\sqrt{p\cdot\sigma}\cdot I\cdot\sqrt{p\cdot\sigma} = p\cdot\sigma = p^0 I - \boldsymbol{p}\cdot\boldsymbol{\sigma}$。
    \item 右下块：$\sqrt{p\cdot\sigma}\cdot I\cdot\sqrt{p\cdot\bar{\sigma}} = mI$。
\end{itemize}
因此：
\begin{equation}
    \sum_\sigma u(\boldsymbol{p},\sigma)\bar{u}(\boldsymbol{p},\sigma) = \begin{pmatrix} mI & p\cdot\bar{\sigma} \\ p\cdot\sigma & mI \end{pmatrix} = \gamma^\mu p_\mu + m \equiv \slashed{p}+m
    \label{eq:pol_u}
\end{equation}
对于反粒子旋量 $v$，由于下半部分负号导致对角块反号：
\begin{equation}
   \sum_\sigma v(\boldsymbol{p},\sigma)\bar{v}(\boldsymbol{p},\sigma) = \slashed{p}-m
    \label{eq:pol_v}
\end{equation}


